\chapter{Implementation}\label{chap:implementation}
% The Engineering Gritty Details: How did you handle code generation for floating-point operations?
% Type Typing/Checking: How does your compiler catch type mismatches (e.g., trying to add an integer to a float if implicit casting isn't supported)?
% Runtime/Backend Integration: If there is a runtime system or specific assembly/IR emitted (like x86 SSE registers or LLVM float primitives), explain how you lowered the AST nodes to it.

The integration of floating-point operations requires a full-vertical extension across the compiler pipeline. This chapter details the engineering semantics, type systems, and lowering mechanisms required to support 64-bit double-precision floats within a base 32-bit architecture.

\section{Interpreter Extension and Phase-Driven Design}

Before implementing the concrete compiler backend, the language's dynamic semantics were established by extending the reference interpreter. Because the interpreter defines the ground-truth behavioral specification of the language, this execution layer must be finalized first to ensure semantic parity with the compiler's eventual target output.

To establish clear verification metrics, a differential test suite containing half a dozen Python programs was constructed. These test cases duplicated existing integer-based control flow, arithmetic expressions, and native built-in functions (such as \texttt{print()}), substituting integer literals with floating-point values. The immediate implementation goal was to allow the interpreter to execute this test suite natively without throwing exceptions, matching Python's runtime behavior exactly.

Extending the interpreter required systematic modifications across three core subsystems of the pipeline: the abstract syntax tree (AST) parser, the static type-checker, and the evaluation semantics engine.

\subsection{Syntactic Extensions (The Parser)}
Extending the frontend parser was highly deterministic because the host language's \texttt{ast} module natively recognizes floating-point tokens. The internal AST grammar definition was upgraded to allow literal expression nodes (\texttt{EConst}) to hold float primitives alongside integers. Additionally, the type-mapping routines were extended to recognize and propagate explicit variable type annotations, allowing patterns such as \texttt{a : float = 3.14} to pass successfully through the front-end token stream.

\subsection{Static Semantics and the 64-Bit Precision Decision}
The primary design challenge emerged during the static analysis phase within the type-checking subsystem. Because the compiler targets a 32-bit RISC-V architecture (RV32), implementing 32-bit single-precision floats (the RISC-V "F" extension) would have offered significant implementation expediency; 32-bit floats map cleanly to standard 32-bit stack slots and integer register pairs. 

However, architectural integrity took precedence over implementation simplicity. Python natively treats all floating-point numbers as 64-bit doubles. Forcing a static 32-bit truncation at the type-checker level would introduce severe semantic drift from standard Python behavior, causing silent precision loss on common practical decimals like \texttt{0.1} or \texttt{3.14}. 

A hybrid architecture—selectively flagging exact base-2 powers (e.g., $0.5 = 2^{-1}$) as 32-bit floats and fallback decimals as 64-bit doubles—was also explored and ultimately rejected. Python does not natively support or expose mixed-precision literals; forcing such an abstraction would create an un-Pythonic dialect. Furthermore, it would introduce significant complexity into the backend code generator, requiring exhaustive permutations of type-coercion routines and hardware conversion instructions for cross-type arithmetic operations. To preserve strict language fidelity, a unified 64-bit double-precision model was selected, targeting the RISC-V Double-Precision ("D") extension. Consequently, the type validator was designed to enforce a single constraint: verifying that literals do not exceed the 64-bit IEEE 754 infinity threshold.

\subsection{Type-System Generalization}
Beyond validating literals, the type-checker's expression evaluation routines needed a clean refactoring. Originally, the compiler assumed almost everything was an integer, evaluating binary operators and functions like \texttt{print()} by checking them directly against a hardcoded \texttt{TInt()} type.
To support floating-point numbers without copying and pasting the same validation logic everywhere, this rigid approach was replaced with a more flexible system. A new helper function, \texttt{check\_types\_supported}, was introduced to validate expressions against a list of acceptable types instead of just one. For this first iteration, operands must still match exactly—meaning mixed-type operations like \texttt{float + int} are blocked for now. However, this refactor ensures that expressions evaluate to and return their correct type primitive (like \texttt{TFloat()}) rather than defaulting to an integer, laying the clean groundwork for future cross-type arithmetic.

\subsection{Interpreter Semantics}
Extending the runtime behavior in the interpreter was the most straightforward step because the interpreter itself is written in Python. Since the host language already supports 64-bit floating-point math natively.

In \texttt{semantics.py}, the evaluation rules for binary expressions were simply updated to accept float values alongside integers. When the interpreter encounters an operation like addition, it handles the AST nodes by passing the raw values directly to Python's native operators. By leveraging the host language's execution environment this way, the interpreter was immediately able to run the floating-point test suite with perfect accuracy.
\section{Extending the compiler to take literals through the whole pipeline}

\begin{itemize}
    \item pass\_1\_2\_shrink: Had to add a 64-bit flag if constant is flow so down the line its exact ieee structure isnt lost.
    \item pass\_1\_2\_shrink: Else the pass didnt have to be adjusted for floats (maybe only the ast adding float)
\end{itemize}

\section{Heap Allocation for Floats}

The decision to box every float carried two practical consequences for the allocation pass. From this pass onwards, a float is represented as a pointer to its box, so any operation that consumes a float must first unbox it to recover the underlying value, and any operation that produces a float must box the result again before it can be stored back on the heap. The difficulty is that the allocation pass must know, for each operand, whether it is a float; otherwise it would apply the arithmetic to the box pointer rather than to the value it points to. The type information needed to make this decision was computed earlier, by the type checker, so the task was to make that information available at this point in the pipeline and to consult it for each operand.

The first step was a function that determines whether a given expression yields a float. Whenever it does, the uniform boxing design dictates that the operand must be unboxed before the operation and, where the operation itself yields a float, boxed again afterwards.

A complication arose for values held inside tuples (and for closures, functions, and objects, which are all represented as tuples at this stage). The type checker does verify the type of every element of a tuple, but that per-element information was not preserved beyond type checking, so by the time the allocation pass ran there was no way to recover whether the element at a given position was a float. The expression's shape alone is sufficient for literals and for nested arithmetic, but not for accesses into a tuple or for plain variables, where the type is not visible in the expression itself. To resolve this, I returned to the type checker and preserved the type information it had already derived.

\subsection{A Dictionary to Preserve Type Information}

The result was a dictionary that, for each such structure, uses its identifier as the key and stores its fields or elements together with their corresponding types as the value. This dictionary is returned from the type checker to the main compiler driver, stored there, and passed on to every pass that requires it.

Up to this point in the pipeline, two passes consume the dictionary. The first is uniquify, which renames variables and must therefore update the dictionary's keys so that the recorded types remain reachable under the new names. The second is the heap allocator, for which the dictionary was originally introduced, where it is used to decide whether a given element is a float.

\subsection{Handling Floats inside the Heap Allocator}

When one or more operands of an operation are floats, the allocator emits a block that first extracts each float value from its box, then performs the operation on the unboxed values, and finally, where the result is itself a float, stores it back into a freshly allocated box. When none of the operands are floats, the operation is emitted unchanged, exactly as before.
\section{Next section}

To be implemented !
